\documentclass[11pt,a4paper]{article}

% --- 极简中文支持（pdfLaTeX 直接兼容，无需切换编译器）---
\usepackage{CJKutf8} % pdfLaTeX 下最稳定的中文支持，无需 XeLaTeX

% --- 基础宏包（清理冗余，确保无冲突）---
\usepackage{amsmath, amsfonts, amssymb} % 数学公式（移除重复的 amsmath）
\usepackage{graphicx} % 插图（已注释调用）
\usepackage{geometry} % 页面布局
\usepackage{listings} % 代码块
\usepackage{xcolor} % 颜色
\usepackage{hyperref} % 超链接
\usepackage{booktabs} % 表格
\usepackage{caption}
\usepackage{float}
\usepackage{subcaption} % 子图支持
\usepackage{multirow} % 表格多行
\usepackage{algorithm} % 算法环境
\usepackage{algorithmic} % 算法伪代码

% 页面设置
\geometry{left=2.5cm, right=2.5cm, top=3cm, bottom=3cm}

% 代码风格设置
\lstset{
    basicstyle=\ttfamily\small,
    keywordstyle=\color{blue},
    commentstyle=\color{green!50!black},
    stringstyle=\color{red},
    frame=single,
    breaklines=true,
    captionpos=b
}

% 标题信息
\title{ruzino曲面参数化实验报告}
\author{宁尚哲 \quad PB24000099}
\date{\today}

\begin{document}
% 开启中文环境（整个文档生效）
\begin{CJK*}{UTF8}{gbsn}

\maketitle

% =============================================================================
% 摘要
% =============================================================================

% =============================================================================
% 1. 引言
% =============================================================================

%实现方案 
\section{数学原理}
%1. 固定边界求解极小曲面
%2. 修改边界条件得到平面参数化
%3. 尝试比较不同的权重

\section{代码实现}  
%代码实现，分两部分：
%需要合适位置讲解半边、3D网络结构数据结构的实现
%param实现
%1. 固定边界求解极小曲面
%pic：节点连接图、balls效果图
%3. 尝试比较不同的权重具体的实现

%boundary map实现
%需要讲解圆形、方形平面映射
%2. 修改边界条件得到平面参数化
%pic：圆形、方形参数化图

\section{纹理可视化样例}
%section4纹理可视化部分
%pic：5个实例模型的可视化
%pic：同模型使用不同权重参数化贴图可视化



\end{CJK*} 
\end{document}